\documentclass[11pt,a4paper]{article}

\usepackage[utf8]{inputenc}
\usepackage[T1]{fontenc}
\usepackage{lmodern}
\usepackage[margin=1in]{geometry}
\usepackage{amsmath,amssymb}
\usepackage{enumitem}
\usepackage{booktabs}
\usepackage{fancyhdr}
\usepackage{xcolor}
\usepackage{mdframed}

\definecolor{solngray}{gray}{0.95}
\newmdenv[backgroundcolor=solngray,linewidth=0.5pt,linecolor=gray,innertopmargin=8pt,innerbottommargin=8pt]{solnbox}

\pagestyle{fancy}
\fancyhf{}
\fancyhead[L]{\textbf{Discrete Structures I}}
\fancyhead[R]{\textbf{Answer Key - Exam D}}
\fancyfoot[C]{\thepage}
\renewcommand{\headrulewidth}{0.5pt}

\begin{document}

\begin{center}
{\LARGE\textbf{ANSWER KEY}}\\[5pt]
{\Large Prelim Practice Exam --- Version D}\\[10pt]
\rule{0.8\textwidth}{0.5pt}
\end{center}

\vspace{15pt}

%==============================================================================
% PART I: MULTIPLE CHOICE
%==============================================================================
\section*{PART I: Multiple Choice (10 points)}

\textbf{1. Answer: (C)}

\begin{solnbox}
$x^2 = 4$ is NOT a proposition because it contains a free variable $x$. Without knowing the value of $x$, we cannot determine if it's true or false.

The other options ARE propositions:
\begin{itemize}[noitemsep]
    \item (A) ``The Earth is flat'' - False, but still a declarative statement
    \item (B) $3 + 4 = 8$ - False, but has a definite truth value
    \item (D) Goldbach's Conjecture - unknown if true or false, but it IS either one or the other
\end{itemize}
\end{solnbox}

\vspace{10pt}

\textbf{2. Answer: (B)}

\begin{solnbox}
The equivalence $p \to q \equiv \neg p \lor q$ is called the \textbf{Implication Law}.

This is crucial for converting conditionals to disjunctions and vice versa. It's often the first step in proving equivalences involving implications.
\end{solnbox}

\vspace{10pt}

\textbf{3. Answer: (A) True}

\begin{solnbox}
Domain: positive reals greater than 1, i.e., $x > 1$.

For any $x > 1$: $x^2 = x \cdot x > x \cdot 1 = x$.

So $x^2 > x$ is true for all elements in this domain. The statement is \textbf{TRUE}.
\end{solnbox}

\vspace{10pt}

\textbf{4. Answer: (C)}

\begin{solnbox}
Negating $\forall x \exists y\, P(x, y)$:

Push the negation inward, flipping each quantifier:
\begin{align*}
\neg \forall x \exists y\, P(x, y) &\equiv \exists x\, \neg \exists y\, P(x, y) \\
&\equiv \exists x \forall y\, \neg P(x, y)
\end{align*}

Remember: negation flips $\forall \leftrightarrow \exists$ and eventually reaches the predicate.
\end{solnbox}

\vspace{10pt}

\textbf{5. Answer: (C)}

\begin{solnbox}
$p \to q$, $q \to r$, therefore $p \to r$ is \textbf{Hypothetical Syllogism}.

It chains implications together. Also called ``transitivity of implication.''
\end{solnbox}

\newpage

%==============================================================================
% PART II: SHORT ANSWER
%==============================================================================
\section*{PART II: Short Answer (20 points)}

\textbf{Question 1.} (5 points) Truth table for $\neg(p \to q) \lor (q \to p)$

\begin{solnbox}
\textbf{Key concepts:}
\begin{itemize}[noitemsep]
    \item $p \to q$ is false only when $p = T$ and $q = F$
    \item $\neg(p \to q)$ is true only when $p = T$ and $q = F$
\end{itemize}

\begin{center}
\begin{tabular}{cc|c|c|c|c}
\toprule
$p$ & $q$ & $p \to q$ & $\neg(p \to q)$ & $q \to p$ & $\neg(p \to q) \lor (q \to p)$ \\
\midrule
T & T & T & F & T & T \\
T & F & F & T & T & T \\
F & T & T & F & F & F \\
F & F & T & F & T & T \\
\bottomrule
\end{tabular}
\end{center}

\textbf{Row 1 detail:} $p \to q = T \to T = T$, so $\neg(T) = F$. $q \to p = T \to T = T$. $F \lor T = T$.

\textbf{Row 3 detail:} $p \to q = F \to T = T$, so $\neg(T) = F$. $q \to p = T \to F = F$. $F \lor F = F$.

This is a contingency (has both T and F).
\end{solnbox}

\vspace{15pt}

\textbf{Question 2.} (5 points) Translation

\begin{solnbox}
``You cannot board the plane unless you have a valid ticket and a passport.''

\textbf{Key phrase:} ``$A$ unless $B$'' can be translated as:
\begin{itemize}[noitemsep]
    \item $\neg B \to \neg A$ (if not B, then not A), OR equivalently
    \item $A \to B$ (if A, then B)
\end{itemize}

The phrase ``cannot [do X] unless [Y]'' means [Y] is necessary for [X].

So: ``You can board'' requires ``(ticket AND passport)''

$$\boxed{b \to (t \land p)}$$

\textbf{Alternative interpretation} (equivalent): $\neg(t \land p) \to \neg b$

Both say: Without both ticket and passport, no boarding.
\end{solnbox}

\vspace{15pt}

\textbf{Question 3.} (4 points)

\begin{solnbox}
\textbf{(i) $\exists x \forall y\, (xy = 0)$}

``There exists a real number $x$ such that for ALL reals $y$, $xy = 0$.''

Yes! $x = 0$ works. $0 \cdot y = 0$ for all $y$.

\textbf{Answer: TRUE}

\rule{\linewidth}{0.3pt}

\textbf{(ii) $\forall x \forall y\, (x + y = y + x)$}

``For all reals $x$ and $y$, $x + y = y + x$.''

This is the \textbf{commutative property of addition}, which holds for all real numbers.

\textbf{Answer: TRUE}
\end{solnbox}

\vspace{15pt}

\textbf{Question 4.} (6 points)

\begin{solnbox}
\textbf{(i) Every student has worked with at least one other student.}

``For every student $x$, there exists a student $y$ (where $y \neq x$) such that $x$ has worked with $y$.''

$$\boxed{\forall x \exists y\, (x \neq y \land W(x, y))}$$

The $x \neq y$ ensures we're talking about ``another'' student, not themselves.

\rule{\linewidth}{0.3pt}

\textbf{(ii) There is a student who has worked with everyone.}

$$\boxed{\exists x \forall y\, W(x, y)}$$

(If we want to exclude self: $\exists x \forall y\, (x \neq y \to W(x, y))$)

\rule{\linewidth}{0.3pt}

\textbf{(iii) Some student has never worked with anyone.}

``There exists a student $x$ such that for all students $y$, $x$ has NOT worked with $y$.''

$$\boxed{\exists x \forall y\, \neg W(x, y)}$$

(If ``anyone else'': $\exists x \forall y\, (x \neq y \to \neg W(x, y))$)
\end{solnbox}

\newpage

%==============================================================================
% PART III: PROBLEM SOLVING
%==============================================================================
\section*{PART III: Problem Solving (25 points)}

\textbf{Question 5.} (6 points)

\begin{solnbox}
\textbf{(i)} Structure: Universal statement about even numbers, plus ``6 is even'', conclude ``36 is even.''

$\forall n\, (E(n) \to E(n^2))$, $E(6)$, therefore $E(36)$

Uses: \textbf{Universal Instantiation} (to get $E(6) \to E(36)$) then \textbf{Modus Ponens}

\rule{\linewidth}{0.3pt}

\textbf{(ii)} $p \lor q$, $\neg q$, therefore $p$

This is \textbf{Disjunctive Syllogism}.

Process of elimination: one option is ruled out, so the other must be true.

\rule{\linewidth}{0.3pt}

\textbf{(iii)} $p \to q$, $\neg q$, therefore $\neg p$

This is \textbf{Modus Tollens}.

If the consequence (pizza) didn't happen, then the cause (Friday) didn't happen either.
\end{solnbox}

\vspace{15pt}

\textbf{Question 6.} (7 points)

\begin{solnbox}
Prove: $\neg(p \lor q) \lor (\neg p \land q) \equiv \neg p$

\begin{align*}
&\neg(p \lor q) \lor (\neg p \land q) \\
&\equiv (\neg p \land \neg q) \lor (\neg p \land q) && \text{[De Morgan's Law]} \\
&\equiv \neg p \land (\neg q \lor q) && \text{[Distributive Law (factor out $\neg p$)]} \\
&\equiv \neg p \land T && \text{[Negation Law: } \neg q \lor q = T\text{]} \\
&\equiv \neg p && \text{[Identity Law: } X \land T = X\text{]}
\end{align*}

$$\boxed{\neg(p \lor q) \lor (\neg p \land q) \equiv \neg p}$$

\textbf{Laws used:}
\begin{enumerate}[noitemsep]
    \item De Morgan's Law: $\neg(p \lor q) = \neg p \land \neg q$
    \item Distributive Law: $(A \land B) \lor (A \land C) = A \land (B \lor C)$
    \item Negation Law: $\neg q \lor q = T$
    \item Identity Law: $X \land T = X$
\end{enumerate}
\end{solnbox}

\vspace{15pt}

\textbf{Question 7.} (4 points)

\begin{solnbox}
Original: $\forall x \forall y\, (P(x) \land Q(y) \to R(x, y))$

\textbf{Step 1:} Negate, pushing through quantifiers.
\begin{align*}
&\neg \forall x \forall y\, (P(x) \land Q(y) \to R(x, y)) \\
&\equiv \exists x\, \neg \forall y\, (P(x) \land Q(y) \to R(x, y)) \\
&\equiv \exists x \exists y\, \neg(P(x) \land Q(y) \to R(x, y))
\end{align*}

\textbf{Step 2:} Simplify the negated implication.

$\neg(A \to B) \equiv A \land \neg B$

So: $\neg(P(x) \land Q(y) \to R(x, y)) \equiv (P(x) \land Q(y)) \land \neg R(x, y)$

$$\boxed{\exists x \exists y\, (P(x) \land Q(y) \land \neg R(x, y))}$$
\end{solnbox}

\vspace{15pt}

\textbf{Question 8.} (8 points)

\begin{solnbox}
\textbf{Premises:}
\begin{enumerate}[noitemsep]
    \item $p \to (q \to r)$
    \item $p$
    \item $q$
\end{enumerate}

\textbf{Goal:} Prove $r$

\begin{center}
\begin{tabular}{c|l|l}
\toprule
\textbf{Step} & \textbf{Statement} & \textbf{Reason} \\
\midrule
1 & $p \to (q \to r)$ & Premise \\
2 & $p$ & Premise \\
3 & $q \to r$ & Modus Ponens (1, 2) \\
4 & $q$ & Premise \\
5 & $r$ & Modus Ponens (3, 4) \\
\bottomrule
\end{tabular}
\end{center}

$$\boxed{\therefore r}$$

\textbf{Explanation:}
\begin{itemize}[noitemsep]
    \item Step 3: Since $p$ is true and $p \to (q \to r)$ holds, by Modus Ponens we get $q \to r$.
    \item Step 5: Since $q$ is true and $q \to r$ holds, by Modus Ponens we get $r$.
\end{itemize}

This demonstrates a ``nested conditional'' being unwrapped in stages.
\end{solnbox}

\vspace{20pt}

\begin{center}
\rule{0.5\textwidth}{0.5pt}\\[5pt]
\textbf{END OF ANSWER KEY}
\end{center}

\end{document}
